\documentclass[a4paper,10pt,twocolumn,landscape]{article}
	\usepackage[ansinew]{inputenc}
	\usepackage{ae}
	\usepackage[brazil]{babel}
	\usepackage{xcolor}
	\usepackage{calc}
	\usepackage{wallpaper}
	\usepackage[margin=2cm]{geometry}
	\usepackage[squaren]{SIunits}
	\usepackage{amsmath}
	\usepackage{paralist}
	\usepackage{lipsum}
	\usepackage{fourier}
	\usepackage{tikz}
	\usepackage{fancyhdr}
	
	\usetikzlibrary{calc}
	
	% Minhas definições
	\newcounter{exercicio}
\newenvironment{exercicio}{%
	\refstepcounter{exercicio}%
	\penalty-200
	\noindent\colorbox{blue!60!black}{\makebox[\columnwidth-\fboxsep*2][c]{\textbf{\color{white}Exercício~\theexercicio}}}\smallskip
}{\par\medskip}
\newcommand{\foreign}[1]{\textsl{#1}}
\newcommand{\digite}[1]{{\fontfamily{cmss}\fontseries{bx}\selectfont#1}}	
\newcommand{\cs}[1]{{\normalfont\textbackslash\color{blue!50!black}#1}}
\newcommand{\pkg}[1]{{\normalfont\sffamily\color{orange}#1}}
\newcommand{\env}[1]{{\normalfont\sffamily\color{green!50!black}#1}}
\let\comando=\cs
\let\package=\pkg
\let\ambiente=\env

	% Papel de parede
	\CenterWallPaper{1}{../figuras/leao-pensador-wattermark.png}

	% Configuração extra: espaço entre as colunas.
	\setlength{\columnsep}{2cm}
	

	
	\newcommand{\cs}[1]{{\normalfont\ttfamily\textbackslash\color{blue}#1}}
	\newcommand{\pkg}[1]{{\normalfont\sffamily\color{orange}#1}}
	
	\newsavebox{\mystarbox}
	\savebox{\mystarbox}{\color{red}$\star$}
	\newcommand{\mystar}{\usebox{\mystarbox}}
	
	\settowidth{\parindent}{\usebox{\mystarbox}}
				
\begin{document}

	\begin{center}\includegraphics[width=0.8\columnwidth]{../figuras/LogotipoCursoLaTeX_v2}\end{center}

	{\noindent\Large Nome: \hrulefill}

	\bigskip

	Na lista de exercícios abaixo, aqueles marcados com \mystar\ devem ser 
	respondidos \emph{nesta} folha. O restante deve ser entregue num arquivo
	\texttt{tex} através do link ``tarefa individual: caixas'' no tópico 8 do
	moodle.

	\bigskip

%------------------------------------------------------------------------------
	\begin{exercicio}
		Qual é a diferença entre os comandos \cs{mbox} e
		\cs{fbox}?		
	\end{exercicio}
%------------------------------------------------------------------------------
	\begin{exercicio}
		O que significa o ``m'' em \cs{mbox}? E o ``f'' em
		\cs{fbox}?
	\end{exercicio}	
%------------------------------------------------------------------------------
	\begin{exercicio}
		Reproduza os resultados abaixo utilizando o comando \cs{framebox}:
		
		\medskip
		
		\begin{center}
			\framebox{Eu sou uma caixa!}\par
			\framebox[5cm][c]{Eu sou uma caixa!}\par
			\framebox[5cm][l]{Eu sou uma caixa!}\par
			\framebox[5cm][r]{Eu sou uma caixa!}\par
			\framebox[1cm][c]{Eu sou uma caixa!}\par
			\framebox[1cm][l]{Eu sou uma caixa!}\par
			\framebox[1cm][r]{Eu sou uma caixa!}
		\end{center}
	\end{exercicio}
%------------------------------------------------------------------------------
	\begin{exercicio}
		Coloque algum texto antes e depois de cada uma das caixas do exercício
		anterior. O que acontece de ``estranho'' nas três últimas caixas? Por quê?
	\end{exercicio}
%------------------------------------------------------------------------------
	\begin{exercicio}
		Reproduza a caixa com dupla moldura:
		
		\medskip
		
		\begin{center}
			\fbox{\fbox{Eu sou uma caixa!}}
		\end{center}
	\end{exercicio}
%------------------------------------------------------------------------------
	\begin{exercicio}
		\noindent\mystar\ Utilize uma régua para traçar uma reta pela linha-base
		da linha abaixo.
	
		\medskip
	
		\begin{center}
			\resizebox{0.99\columnwidth}{!}{Identifique a linha-base.}
		\end{center}
		
		\medskip
		
		\footnotesize
		\noindent\textbf{extra:} trace uma reta \unit{1}{ex} acima da 
		linha-base\dots\ você não precisa medir nada, apenas olhar.
	\end{exercicio}
%------------------------------------------------------------------------------
	\begin{exercicio}
		\noindent\mystar\ No exercício anterior, esboce as diversas caixas que
		compõem aquela linha e seus respectivos pontos-de-referência.
		
		\medskip
		
		\footnotesize
		\noindent\textbf{obs.:} a seqüência ``fi'' é, na verdade, um caráter só. É
		o que se chama \emph{ligação} ou caracteres-ligados (\foreign{ligatures}) e
		serve para melhorar a aparência de algumas seqüências de letras. Nas fontes
		Computer Moder Roman, exemplos são fi, ff, fl, ffi e ffl, mas em outras
		fontes (como esta: utopia) outras seqüências podem ser definidas.
	\end{exercicio}
%------------------------------------------------------------------------------	
	\begin{exercicio}
		\noindent\mystar\ Complete: o parâmetro \cs{baselineskip} rege a distância 
		entre as \hrulefill\\\null\hrulefill de duas linhas consecutivas.
	\end{exercicio}
%------------------------------------------------------------------------------
	\begin{exercicio}
		\noindent\mystar\ Esboce as caixas que contém as caixas-pretas e seus 
		pontos-de-referência.
		
		\medskip
		
		\Huge
		\noindent\makebox[0pt][l]{\raisebox{1pt}{\scriptsize Linha-base}}%
		\hrulefill\
		\rule[2ex]{1em}{1ex}
		\rule[-0.5ex]{1em}{1ex}	
		\rule[0ex]{1em}{1ex}		
		\rule[-1ex]{1em}{1ex}
		\rule[1ex]{1em}{1ex}
		\rule[-2ex]{1em}{1ex}
		\rule[0.5ex]{1em}{1ex}
		\hrulefill
	\end{exercicio}
%------------------------------------------------------------------------------
	\begin{exercicio}
		\noindent\mystar\ No final do texto abaixo, a palavra ``purus'' foi
		separada pelo \TeX\ no final da linha. Utilizando o que você aprendeu
		sobre caixas, como você poderia evitar que isso acontecesse?
		
		\medskip
	
		\large
		\foreign{%
		Lorem ipsum dolor sit amet, consectetuer adipiscing elit. Ut purus elit,
		vestibulum ut, placerat ac, adipiscing vitae, felis. Curabitur dictum
		gravida mauris.}
		
		\medskip
		
		\noindent\footnotesize\textbf{obs.:} \emph{não} use \cs{hyphenation}.

	\end{exercicio}
%------------------------------------------------------------------------------
	\begin{exercicio}
		Considere o resultado abaixo. Sabendo que existe um comando que reflete uma
		caixa com relação à linha-base (isto é, assuma que colocar o texto de 
		cabeça para baixo não é um problema), \textbf{explique} como você faria
		para construir o resultado abaixo.
		
		\medskip
		
		\begin{center}
		\Huge\makebox[0pt][l]{\raisebox{0.5pt}{Eu sou uma caixa!}}%
		\scalebox{1}[-1]{\textcolor{blue!30}{\raisebox{0.5pt}{Eu sou uma caixa!}}}
		\end{center}
				
	\end{exercicio}
%------------------------------------------------------------------------------
	\begin{exercicio}
		\textbf{DESAFIO:} reproduza o resultado do exercício anterior.
		
		\medskip
		
		\footnotesize
				
		\noindent\textbf{obs.:} a instrução \cs{scalebox}\verb|{1}[-1]{|%
		\cs{textcolor}\verb|{blue!30}{Eu sou uma caixa!}}| produz o texto
		invertido em azul. \textbf{Dica:} utilize \cs{makebox}.		
	\end{exercicio}
%------------------------------------------------------------------------------
	\begin{exercicio}
		\noindent\mystar\ Complete a linha-base da linha abaixo e esboce os
		pontos-de-referência das caixas que englobam cada palavra.
		
		\medskip
	
		\noindent\hrulefill
		\raisebox{1ex}{Reproduza}
		\raisebox{0ex}{esta}
		\raisebox{-1ex}{linha}
		\raisebox{0ex}{de}
		\raisebox{1ex}{texto}
		\raisebox{0ex}{utilizando}
		\raisebox{-1ex}{o}
		\raisebox{0ex}{comando}
		\raisebox{1ex}{\cs{raisebox}}
		\raisebox{0ex}{.}
		\hrulefill
		
		\medskip
		
		\footnotesize
		\textbf{obs.:} além de desenhar as caixas, você deve reproduzir a linha 
		acima e enviá-la ao professor.
	\end{exercicio}
%------------------------------------------------------------------------------
	\begin{exercicio}
		Crie uma linha preta de \unit{1}{pt} de espessura que vai da margem direita
		até a margem esquerda da página, como esta:
		
		\medskip
		
		\noindent\rule{\columnwidth}{1pt}
		
		\medskip
		
		\noindent\textbf{obs.:} não se esqueça de suprimir o avanço do parágrafo
		com \cs{noindent}.
	\end{exercicio}
		
%------------------------------------------------------------------------------
	\begin{exercicio}
		Refaça o exercício anterior, mas desenhe uma linha vermelha:
		
		\medskip
		
		{\color{red}\noindent\rule{\columnwidth}{1pt}}
		
	\end{exercicio}
%------------------------------------------------------------------------------
	\begin{exercicio}
		Reproduza o resultado abaixo utilizando o comando \cs{parbox}.
		
		\medskip
		
		\begin{center}
			\tiny
			\parbox[t]{0.4\textwidth}{\lipsum[1]}\qquad%
			\parbox[t]{0.4\textwidth}{\lipsum[3]}
		\end{center}
		
		\smallskip
		
		\footnotesize
		\noindent\textbf{dicas:} a distância entre as duas colunas (caixas) é
		\cs{quad}. Utilize fontes de tamanho \cs{tiny} e, se quiser, o pacote
		\pkg{lipsum}.
	\end{exercicio}
%------------------------------------------------------------------------------
	\begin{exercicio}
		Refaça o exercício anterior usando mini-páginas.
	\end{exercicio}
%------------------------------------------------------------------------------
	\begin{exercicio}
		O espaço adicional que precede a primeira linha de um parágrafo, também
		conhecido por \textbf{e}ndentação, tem comprimento parametrizado por
		\cs{parindent} (\foreign{paragraph indentation}).
		
		Como você pôde ver nos dois exercícios anteriores, este comprimento é zero
		nas caixas de parágrafo (em qualquer caixa, na verdade). Seu trabalho 
		é modificar isso: introduza alguma endentação no texto dessas caixas.
	\end{exercicio}
%------------------------------------------------------------------------------
	\begin{exercicio}
	Para colocar uma moldura numa expressão matemática em estilo de exibição, 
	você deve antes colocá-la dentro de uma caixa de parágrafo. Faça isso:
	reproduza o resultado abaixo.
			
	\medskip
	
	\noindent\framebox[\columnwidth]{%
		\parbox{\columnwidth}{%
			\[
				\int\limits_0^{2\pi} \sin(nx)\sin(mx)\, dx = \delta_{nm}
			\]}}
	\end{exercicio}
	
\end{document}